\documentclass[a4paper,12pt]{article}
\usepackage[utf8]{inputenc}
\usepackage{geometry}
\geometry{margin=1in}
\usepackage{hyperref}
\usepackage{booktabs}
\usepackage{graphicx}

\title{\textbf{Project Report} \\ Hardware-Accelerated Edge-AI Security System}
\author{\textbf{Roll No: 23f3000570}}
\date{\today}

\begin{document}

\maketitle

\begin{abstract}
This report details the design and deployment of a privacy-first, low-power Edge-AI surveillance platform built on embedded hardware. The system integrates a Raspberry Pi 5 Single Board Computer (SBC) with a Hailo-8L Neural Processing Unit (NPU) via a PCIe Gen3 x1 interface. By offloading deep learning inference to the dedicated NPU, the system achieves real-time multi-stage computer vision at 52.2 FPS within a sub-10 W power envelope, ensuring all processing remains on-device without cloud dependency.
\end{abstract}

\section{Introduction}
Modern smart home security systems typically stream video data to cloud servers due to the high computational requirements of deep learning models, raising significant privacy concerns. This project addresses this limitation through hardware-software co-design, leveraging modern embedded accelerators. The primary objective is to build a standalone, power-efficient electronics platform capable of continuous real-time video inference, secondary verification, and secure network communication entirely at the edge.

\section{Hardware Architecture and Component Selection}
The system is built around a distributed compute architecture, decoupling the main application logic from neural network execution.

\begin{itemize}
    \item \textbf{Main Processing Unit:} A Raspberry Pi 5 SBC featuring a quad-core ARM Cortex-A76 (BCM2712) clocked at 2.8 GHz, equipped with 16GB LPDDR4X RAM. This unit handles video capture, web services, and asynchronous event routing.
    \item \textbf{Hardware Accelerator (NPU):} A Hailo-8L AI HAT+ delivering 13 TOPS of INT8 compute. It interfaces with the host SBC via an M.2 M-key slot utilizing a single-lane PCIe Gen3 link. This enables Direct Memory Access (DMA) for video frames, bypassing the host CPU during forward pass inference.
    \item \textbf{Vision Sensor:} Raspberry Pi Camera Module v3 (Sony IMX708 sensor) connected via the MIPI CSI-2 interface, providing high-resolution raw frames directly to the hardware pipeline.
    \item \textbf{Storage and Power:} Persistent storage is handled by a 1TB SATA SSD over USB 3.0. The entire rig is powered by a high-current 5V/5A USB-C Power Delivery (PD) supply. Thermal management is achieved using the official active cooler.
\end{itemize}

\section{Embedded Software and Inference Pipeline}
The software stack is highly optimized for the chosen hardware:
\begin{itemize}
    \item \textbf{Inference Engine:} The vision pipeline is constructed using GStreamer and Hailo TAPPAS. A custom INT8-quantized YOLOv8s object detection model runs directly on the Hailo-8L NPU. 
    \item \textbf{Multi-stage Processing:} A secondary MobileNetV2 classifier and MiDaS Small depth-estimation model share the NPU's resources through time-multiplexed \texttt{infer} calls.
    \item \textbf{CPU-side Services:} The Cortex-A76 cores handle a FastAPI web server, AES-256-CBC encrypted video storage, secure remote ingress via Cloudflare tunnels, and a local voice assistant daemon. Thread-affinity pinning is utilized to prevent CPU scheduler jitter from interfering with the camera capture pipeline.
\end{itemize}

\section{System Performance and Thermal Characteristics}
The integrated hardware platform was subjected to rigorous empirical evaluation:

\begin{table}[h]
\centering
\caption{Measured Hardware Performance Metrics}
\begin{tabular}{@{}ll@{}}
\toprule
\textbf{Parameter} & \textbf{Measured Value} \\ \midrule
End-to-End Pipeline Throughput & 52.2 FPS (sustained) \\
NPU On-Chip Latency & 18.4 ms \\
Total Board Power Draw (Mean) & $5.89 \pm 1.43$ W (measured via PMIC) \\
Peak Power Draw & 8.58 W \\
NPU Surface Thermal Rise & $+3.8^\circ$C over ambient idle \\
Host CPU Median Load & 48.2\% \\ \bottomrule
\end{tabular}
\end{table}

The NPU effectively isolates the high-density matrix multiplication workloads, resulting in a system that maintains real-time operation well above the 30 FPS surveillance threshold. The measured sub-10 W total power envelope demonstrates the viability of the system for continuous 24/7 deployment.

\section{Conclusion}
This project successfully demonstrates the integration of high-performance neural accelerators with affordable Single Board Computers. By carefully managing PCIe bandwidth, thread affinity, and INT8 model quantization, the resulting electronics platform delivers a robust, secure, and fully localized AI surveillance solution. It serves as a practical reference architecture for privacy-first, low-power edge computing.

\end{document}
